\section{Design Space}
\label{sec:design}

All schemes authenticate the same canonical claim string, so only the authenticator and its delivery differ. Fig.~\ref{fig:arch} shows the two post-quantum designs this paper focuses on.

\begin{figure}[t]
\centering
\begin{tikzpicture}[
  font=\footnotesize,
  box/.style={draw, rounded corners=2pt, align=center, minimum height=0.7cm, inner sep=2.5pt, fill=white},
  lab/.style={font=\scriptsize, align=center, fill=white, inner sep=1pt},
  arr/.style={-{Stealth[length=1.8mm]}, thick},
  darr/.style={-{Stealth[length=1.8mm]}, thick, dashed}]
\node[box, minimum width=2.4cm] (iss) at (0,0) {Issuer $I$\\ M2: sign each token\\ EMBER: sign root $R_e$};
\node[box, minimum width=2.2cm] (rs) at (5.0,0) {Root service\\ (EMBER only)};
\node[box, minimum width=2.2cm] (sh) at (5.0,-1.55) {Shield tier};
\node[box, minimum width=2.2cm] (edge) at (5.0,-3.2) {Edge $E_j$\\ $pk$, verified cache};
\node[box, minimum width=2.4cm] (cli) at (0,-3.2) {Client};
\draw[arr] (iss) -- node[lab, above]{$SR_e$} (rs);
\draw[arr] (rs) -- (sh);
\draw[arr] (sh) -- node[lab, left=2pt]{pull on miss\\(fallback)} (edge);
\draw[darr] (rs.east) .. controls +(0.9,-0.5) and +(0.9,0.5) .. node[lab, right]{push} (edge.east);
\draw[arr] (iss) -- node[lab, right=2pt, pos=0.45]{URL token\\$+$ signature or $SR_e$\\(to staple)} (cli);
\draw[arr] (cli) -- node[lab, below=3pt, pos=0.42]{requests with URL token;\\ staple per connection} (edge);
\end{tikzpicture}
\caption{Stapled post-quantum tokens. In M2 the client staples its token's ML-DSA signature on the first request of each connection. In EMBER it staples the signed epoch root $SR_e$ (MTL's ``full'' signature) and the URL carries a Merkle path. Pull from the root service through the shield tier and push are EMBER fallbacks.}
\label{fig:arch}
\end{figure}

\subsection{Per-Token Signatures, Cached or Stapled}
\textbf{B5/B6} put an ML-DSA-44/65 signature in every token. \textbf{B5v} adds an edge cache of verified tokens, keyed by SHA-256 of the token, so verification runs once per (token, edge). The 3.3~KB token is still sent and hashed on every request. \textbf{M2} (stapled per-token ML-DSA-44) separates the signature from the token. The URL token is the claims followed by \texttt{:S=}b64url($\mathit{kid}\,\|\,$16-B nonce). On the first request of each connection, the client sends the 2,420-B signature over (claims, nonce) in a header. We use the connection as the staple unit because a client can observe connections but not edges. The edge verifies the signature, caches $H(\mathit{ref}\,\|\,\text{claims})$ as verified until the token's expiry, and from then on hashes about 138~B per request.

\emph{Recovery.} M2 has no pull fallback. An edge that holds no valid entry for a token answers with 401 and a \texttt{Staple-Required} hint, and the client retries the same request with the signature header. This happens after a cache eviction, an edge process restart, or a server that does not share the cache. Our prototype implements the hint as a 403 with \texttt{X-Reject: staple required}. A recovery costs one client--edge round trip (20.98~ms mean in our testbed), a \num{3242}-B HTTP/1.1 header line, and one ML-DSA-44 verification (100.9\us). Because clients staple on every new connection, recovery is needed only when an edge loses an entry for a connection that is still open. The edge state is small. In the prototype an entry is a 16-B key plus an 8-B expiry (24~B, excluding hash-table overhead), and it lives at most until the token expires ($\le\vmax=1$~h). At the busiest simulated edge, new (token, edge) pairs arrive at 165.8/s ($\bar\lambda=2000$/s) and \num{1658}/s (\num{20000}/s). This bounds the live set at 0.60 and 5.97 million entries, or 14.3 and 143~MB of payload. An edge that retains this state never needs recovery for unexpired tokens. We did not measure how often recovery occurs. M2 is the $N=1$ special case of EMBER, and it is also B5v without the per-request resend. Its security is the EUF-CMA security of ML-DSA plus the collision resistance of the cache key.

\subsection{EMBER: Pooled Epoch-Merkle Batching}
\label{sec:ember}
\emph{Notation.} $H$ is SHA-256 truncated to $n\in\{16,32\}$ bytes. The tweakable hash is $\Th(T,x)=\mathrm{trunc}_n(H(T\|x))$ with a fixed-length 26-B tweak $T=\texttt{"EMBERv1\textbackslash0"}\|\mathit{kid}\|e\|\mathit{type}\|\ell\|j\|\mathit{tenant}$, so every node position of every epoch is a distinct function~\cite{sphincsplus2019}. Leaf $i$ of epoch $e$ is $L_i=\Th(T_{\textsc{leaf}}(e,i),\,r_i\|\mathrm{enc}(c_i))$, where $r_i$ is a 128-bit salt that the issuer draws \emph{after} it receives the claims $c_i$. Inner nodes are $\Th(T_{\textsc{node}}(e,\ell,j),\,\text{left}\|\text{right})$, and the tree has the RFC~9162 shape~\cite{rfc9162}, with depth $d=\lceil\log_2N_e\rceil$. The signed epoch root is $SR_e=(\mathit{kid},e,T_{\text{start}},\Delta_e,N_e,R_e,\sigma_e)$, where $\sigma_e=\Sig(sk,\texttt{"EMBER-root"}\|\mathit{kid}\|e\|T_{\text{start}}\|\Delta_e\|N_e\|R_e)$ is ML-DSA-44 with the FIPS~204 context string \texttt{EMBER-root}. In the hybrid variant, $\sigma_e$ is an Ed25519 signature and an ML-DSA-44 signature over the same message, and both must verify.

\emph{Token.} The URL token is \texttt{URLPrefix=\ldots:Expires=\ldots:} \texttt{Tenant=\ldots:E=}\textit{blob}, where the blob is the base64url encoding of version, algorithm, $\mathit{kid}$, $e$, $i$, $r_i$ and the path, $30+nd$ raw bytes in all. With $c_{\text{enc}}$ bytes of encoded claims,
\begin{equation}
S_{\text{tok}}=c_{\text{enc}}+3+\lceil 4(30+nd)/3\rceil .
\label{eq:size}
\end{equation}
Our encoder implements Eq.~\eqref{eq:size}. With a 60-character path ($c_{\text{enc}}=118$) it yields 161, 268, 332 and 460~B for $N=1$, $2^5$, $2^8$ and $2^{14}$, against 148~B for M2.

\emph{Issuance and pooling.} The issuer collects requests for $\Delta$ seconds, or until a count cap $N_{\max}$ is reached, and then closes the epoch. Closing builds the tree ($2N-1$ hashes), signs the root once, and returns to each client its token and $SR_e$. Before the epoch or index counters would wrap, the issuer rotates $\mathit{kid}$. The unit of batching is the \emph{pool}, meaning the set of tenants that share a tree. Batching saves signatures only when a pool receives many tokens per $\Delta$. With $P$ pools the issuer makes at most $P/\Delta$ signatures per second, independent of the token rate. Pooling lets one issuer key mint tokens for every tenant in the pool. It is therefore appropriate only when the tenants already delegate issuance to the CDN or belong to one trust domain. Tenant is part of the claims and of the leaf tweak, so leaves cannot be replayed across tenants.

\emph{Verification} is given in Algorithm~\ref{alg:verify}. Revocation uses short $\vmax$ (1~h by default), signed lists of revoked $(e,i)$ pairs at 8~B each, or withdrawal of a whole epoch root.

\begin{algorithm}[t]
\caption{Edge verification of an EMBER token}
\label{alg:verify}
\small
\begin{algorithmic}[1]
\Procedure{Verify}{$\mathit{tok},\mathit{url},\mathit{now}$}
  \State parse claims $c$ and $(\mathit{kid},e,i,r,\text{path})$; reject if malformed or $\mathit{kid}$ unknown
  \State reject if $\mathit{now}>c.\mathit{Exp}+\tau$ or $c.\mathit{Exp}>T_{\text{start}}(e)+\Delta_e+\vmax+\tau$
  \State reject if URL prefix or IP binding fails, $(e,i)$ revoked, or $i\ge N_e$
  \State $R\gets\mathit{RootCache}[\mathit{kid},e]$
  \If{$R=\bot$}
    \If{a stapled $SR_e$ is present} verify $\sigma_e$ (per-client rate limit); cache
    \Else{} fetch $SR_e$ via shield (coalesced, negative-cached); verify; cache
    \EndIf
    \State on failure or timeout: \Return reject \Comment{fail closed}
  \EndIf
  \State $x\gets\Th(T_{\textsc{leaf}}(e,i),r\|\mathrm{enc}(c))$
  \For{each sibling $s_\ell$ in path} $x\gets\Th(T_{\textsc{node}}(e,\ell,\cdot),\mathrm{order}(x,s_\ell))$
  \EndFor
  \State \Return $x\stackrel{\text{ct}}{=}R$
\EndProcedure
\end{algorithmic}
\end{algorithm}

\subsection{Root Dissemination}
\label{sec:dissem}
An edge cannot verify an EMBER token until it holds $R_e$. There are four ways to get it there. \textbf{Pull}: the edge fetches $SR_e$ through the shield on a miss, with coalesced fetches. \textbf{Push}: the issuer fans each root out to all edges, at fleet bandwidth $BW_{\text{push}}=T_{\text{active}}\,E\,(|\sigma|+40)/\Delta$ for $T_{\text{active}}$ active pools. \textbf{Period roll-up}: every $m$ epochs one signed bundle lists the $m$ roots. \textbf{Stapling}: the client sends $SR_e$ in a header, either only on the session's first request (\emph{staple-first}, our earlier design) or on the first request of every connection (\emph{staple-all}). In our workload a session opens one connection per edge it uses, so staple-all equals stapling on every new edge contact, which is how the simulator models it. The stapled header is the 2,461-B signed root, \num{3298}~B as an HTTP/1.1 header line. With staple-all, verification does not depend on the root service, which matters because edges fail closed. During an outage of the root service or the shield tier, pull-only edges reject every token of an epoch they have not cached, while stapled tokens still verify. M2 and staple-all EMBER instead depend on the edge retaining verified entries, or on the recovery retry. For Poisson token arrivals in pull mode, the expected number of root misses per token at edges with traffic shares $w_v$ is
\begin{equation}
M(\Delta)=\frac{1}{\lambda\Delta}\sum_{v=1}^{E}\bigl(1-e^{-\lambda\Delta w_v}\bigr).
\label{eq:miss}
\end{equation}
The added delay per token is $D=\Delta/2+t_{\text{epoch}}+M(\Delta)F_{\text{miss}}$, where $F_{\text{miss}}$ is the root round trip plus $t_v$ (pull), or $t_v$ alone (stapling). When $\lambda\Delta\ll E$, $M\to1$: almost every token is the first at its edge in its epoch, so only stapling or push avoids the round trip.

\subsection{Security}
\label{sec:security}
\begin{theorem}
\label{thm:uf}
Suppose an adversary $\mathcal{A}$ makes $q_h$ queries to $\Th$ and the issuer evaluates $\Th$ at $p\le\sum_e(2N_e-1)$ tree positions. Then there exist $\mathcal{B}_1,\mathcal{B}_2$ with about the same running time such that
$\Adv^{\mathrm{uf\text{-}tok}}_{\ember}(\mathcal{A})\le\Adv^{\mathrm{euf\text{-}cma}}_{\Sig}(\mathcal{B}_1)+\Adv^{\mathrm{sm\text{-}tcr}}_{\Th,p}(\mathcal{B}_2)$.
\end{theorem}
\begin{IEEEproof}[Proof sketch]
Let $R^*$ be the root the edge used. If the issuer never produced $R^*$ for $(\mathit{kid},e^*)$, then the verified $\sigma$ is a forgery, which $\mathcal{B}_1$ outputs. Otherwise $R^*=R_{e^*}$ is honest, and since $i^*<N_{e^*}$ is enforced, an honest leaf $i^*$ exists whose claims differ from $c^*$. The honest and adversarial hash chains start from different inputs and reach the same root. So at some position they have different inputs and equal outputs under the same tweak. The honest input there contains the salt $r_{i^*}$, which was chosen after $\mathcal{A}$ fixed $c_{i^*}$, or a value derived from it. $\mathcal{A}$ therefore hits a fixed honest target, which is an sm-tcr break; distinct tweaks rule out multi-target speed-ups~\cite{sphincsplus2019}. Replays across epochs, tenants or $\mathit{kid}$ change the tweak and fall into the second case.
\end{IEEEproof}
This is a sketch, not a new proof. For $n=32$, plain collision resistance suffices, matching batch signatures~\cite{batchsig2024}. The root and period-bundle signatures use distinct context strings. The hybrid root is an AND-combiner, at least as strong as its stronger component~\cite{bindel2017hybrid}. sm-tcr is defined for a keyed hash family. Our tweak contains $\mathit{kid}$ and the epoch but no random public seed (SLH-DSA's PK.seed). We therefore recommend adding a 16-B random seed per $\mathit{kid}$ to the tweak, carried in $SR_e$. This adds 16~B to the staple and nothing to tokens. We did not measure this variant. \emph{Truncation:} ML-DSA-44 is a category-2 parameter set~\cite{fips204}. With $n=16$, the generic second-preimage bound is about $2^{128}$ classically and $2^{64}$ with Grover. The \emph{token} therefore offers category-1 security, the same generic level as SLH-DSA-128's 16-B hashes~\cite{fips205}. Use $n=32$ for higher categories (+341~B at $N=2^{16}$). The issuer-chosen salt is mandatory, because $\mathcal{A}$ must not control both inputs of a collision.

\emph{Other threats.} Downgrade: once a tenant has switched, edges must reject classical tokens for that tenant. Privacy: $N_e$ in a public root leaks issuance volume, and the pair $(e,i)$ identifies a token, as a per-token signature already does. Padding the tree to $2^d$ hides the volume. It slowed our Python issuer by 12.45\% at $N=1000$, but the worst case is $N=2^k+1$, where the analytic native cost rises by 38.9\% at $N=2^{14}+1$. DoS: tokens naming fake epochs trigger root fetches, and fake staples trigger signature verifications (Section~\ref{sec:res-other}). Replaying a \emph{valid} staple costs nothing at an edge that has the root cached, and costs at most one verification per $(\mathit{kid},e)$ elsewhere. Per-client limits bound both cases.
